% Pursuing Stacks (cotes 134-2 to 134-8): archive cote 134-5, site page 69 (PDF page 70), typescript page 256.
% typescript. Draft compared: ps4.tex:3480-3536.
% AI-transcribed by Claude Opus 5.5 (claude-opus-5-5, adaptive thinking, effort high), Message Batches API, batch ps2 (msgbatch_01NVwr1Yuo8qJf5rykeGBXPK), 2026-11-12; not yet checked by eye.
\scanpage{134-5/69}
\mspage{256}

associating to a notion of weak equivalence in M i.e. to any saturated
subset \xout{$W_a$}
\[ W_a \subset \mathrm{Fl}(M) \quad , \]
a corresponding homotopy structure $h_{W_a}$ ,\xout{,} a homotopy interval structure
as a matter of fact (section 54, p. 131).  Here we are primarily interes-%
ted of course in the case when $W_a$ is associated to a\deleted{n} \inserted{given} asphericity structu
-re $M_{as}$ in M , which we may as well assume to be aspheric, so \deleted{\unsure{xas}} \inserted{as} to
be sure that $W_a$ and $M_{as}$ determine each other.  I recall that the weak
interval structure $h_{W_a}$ \xout{is} can be described by the generating family of
homotopy intervals, consisting of all intervals
\[ \underline{I} = (I,\hand{\delta}_o,\hand{\delta}_1) \]
such that $I \to e_M$ be universally in $W_a$ , i.e. such that for any object
x in M, the projection
\[ x \times I \longrightarrow x \]
be in $W_a$ .  \deleted{In view of} \inserted{For} pinning down further the exact relationship bet-%
ween contractibility structures (which \xout{are} \inserted[typed]{may\xout{x} be viewed as} just special types of homoto-%
py interval structures) and asphericity structures, we are thinking of
course \xout{(in} more specifically of \uline{totally} \uline{aspheric} asphericity structu-%
res, in view of prop.2 c) of the preceding section (p. 248).  It is \deleted{rea-%
sonable to expect} \inserted{immediate} in this case that \xout{xxx} \inserted[typed]{for} an object I of M,\xout{is such that}
$I \to e_M$ is in $UW_a$ (i.e. is universally in $W_a$) \deleted{iff} if (and only if,
of course) I is aspheric, i.e. $I \longrightarrow e_M$ is in $W_a$ .  \deleted{But maybe I'll
leave \xout{this} elucidation of this point for a later moment, in order not
to diverge from \xout{the}} \over{t}{T}he mo\unsure{s\over{e}{t}} relevant questions which come up here, \inserted{now}
\deleted{namely:} \inserted{seem to me the following:}

1)  If $M_{as}$ is generated by a contractibility structure $M_c$ ,  is the
homotopy structure $h_{W_a}$ associated to $M_{as}$  (indeed, a \xout{weak} homotopy in-%
terval structure as recalled above) also the one defined by $M_c$ ,  using
intervals in $M_c$ as a generating family of homotopy intervals\unsure{\over{.}{?}}

2)  Conversely, what extra conditions on a given asphericity struc-%
ture $M_{as}$ on M are needed (besides total asphericity) to insure that the
corresponding homotopy structure $h_{W_a}$ on M \xout{is a} \inserted[typed]{comes from a} contractibility structure \inserted{$M_c$}
(i.e. admits a generating family of homotopy intervals which are \uline{contrac-}%

\margin{69}
